\section{The extension boundary}
\label{sec:scope}

\subsection{What the inspected repository already contains}

The inspected Forge snapshot is not an empty proposal. Its merged article is
111 pages, draws on eighteen preserved submissions, and already contains the
usual candidates for a stronger proof planner. Its canonical status ledger,
provenance appendix, conclusion, and Leant/Djex review were used to delimit this
work~\cite{forge,forge-status,forge-provenance,forge-ideas}. In particular,
induction with generalization, arithmetic witness programs, ideal closure,
finite-algebra covers, observable-space closure, Kripke countermodels,
nonlinear certificates, boundary-safe telescoping, and cyclic descent are
\emph{prerequisites}, not contributions of this article.

The merge also records successful elaboration of sixteen core-only Lean files.
It would therefore be wrong to say that nothing in Forge has ever compiled.
The relevant remaining gap is narrower: the canonical ledger does not claim an
implemented \forge{} tactic or a Lean-verified certificate checker for its
workers~\cite{forge-status}. This article does not close that gap merely by
adding Python experiments.

The non-duplication claim here is scoped to the canonical merged design and
capability inventory at the pinned commit. It is not a claim that every line of
every historical submission has been exhaustively searched, nor that the
algorithms below are new to the literature. The repository itself encourages a
small checked implementation over another repetition of its architecture; the
proposal below keeps that architecture and changes the kinds of obligations its
workers can settle.

\begin{longtable}{@{}L{0.25\linewidth}L{0.32\linewidth}L{0.36\linewidth}@{}}
\toprule
Existing mechanism & What it does not itself provide & Addition in this article \\
\midrule
\endhead
Universal word/observable closure & A controller's choices against an adversary,
with an infinite recurrence condition & Winning positional strategies plus
threshold-local parity ranks \\
Finite-state or finite-algebra closure & Correlated executions with prescribed
probability marginals & Full rational couplings and optimal unsupported-mass
certificates \\
Closure under all transitions & The alternating order ``for every source action,
there exists a matching target action'' & Greatest probabilistic simulation with
replayable elimination evidence \\
Well-founded cyclic descent & A loop that may repeat forever as a graph path,
but terminates almost surely & Bellman drift certificates for expected time,
or a reachable controlled trap \\
Typed witness synthesis & Behavioral correctness of a strategy or a coupled
transition table & Source-owned strategy tables and probabilistic witnesses,
with separate semantic proof obligations \\
\bottomrule
\end{longtable}

\subsection{Why these are related capabilities}

A useful way to view the extension is to distinguish three transition
operators. For a predicate $P$, a demonic step requires every successor to
satisfy $P$, whereas an angelic step requires a selectable successor. For a
nonnegative numerical potential $V$, a probabilistic step transports $V$ by an
expectation. The same proof-planning mistake can occur in all three settings:
the operator is silently changed to an easier one.

Replacing $\forall a\,\exists b$ by $\exists b\,\forall a$ loses valid strategy
witnesses; replacing it by $\exists a\,\exists b$ accepts invalid ones. Replacing
an expectation by pointwise decrease loses randomized termination proofs.
Replacing a coupling by an independent product loses valid relational proofs.
Treating a finite occurrence of a good priority as recurring accepts a false
liveness claim. Each worker below carries enough information for its checker
to distinguish precisely these cases.

The shared design principle is consequently more specific than ``use another
solver'': \emph{reify the transition quantifier, propose a finite witness, and
check its local obligations under the original operator}. Probability mass,
choice ownership, terminal semantics, and the observation interface are part of
the problem, not conveniences for a solver to reconstruct.

\subsection{A concrete target workload}

Four kinds of Lean goals motivate the workers. The following is schematic
notation, not syntax installed by the accompanying package:
\begin{lstlisting}
-- A controller exists against every legal environment strategy.
exists controller, forall environment,
  parityWins (run arena controller environment start)

-- Two probabilistic computations admit a relational joint execution.
exists joint, marginals joint = (law left, law right)
  and probability joint (not relation) <= epsilon

-- Every implementation action admits a matching specification action.
probabilisticSimulation source specification initialPair

-- Every scheduler has bounded expected termination time.
forall scheduler, expectedSteps program scheduler start <= budget
\end{lstlisting}

These are not four names for a finite reachability computation. The first is an
infinite-game statement. The second has a joint-distribution witness. The third
has alternating choices before sampling. The fourth quantifies over schedulers
and estimates a numerical random variable. The finite input fragments let each
worker return a small certificate even though the semantic claim concerns
arbitrarily long executions.

\subsection{What was actually built}

The archive contains a new package, \code{forge_ap}, rather than a modified copy
of Forge. Its search module implements recursive Zielonka solving, rational
max-flow, a greatest-simulation elimination loop, and exact policy iteration.
Its checker module performs only local graph, sum, and inequality checks; it
imports no search algorithm. Separate tiny-model oracles use strategy
enumeration, Hall-subset enumeration, exhaustive subrelations, and determinant
expansion. This is stronger separation than replacing a solver entry point with
an exception while retaining the solver's arithmetic assembler, although the
implementations still share representations and Python's arithmetic runtime.

No compiler was available for the Lean part of this execution. Accordingly,
\code{lean/} contains small, explicitly uncompiled lowering specimens; the
experiment ledger records \status{NOT\_RUN}. The mathematical proofs in the
article explain what a future Lean checker must establish. They do not turn
Python execution into kernel acceptance.

\subsection{Pinned context and scope of external tools}

The inspected repository revisions are recorded in Appendix~\ref{app:artifacts}.
Forge's pinned Lean baseline is v4.34.0 with Mathlib revision
\texttt{1cf325a0cf67aca2b04d76b5380ff6a9e410aefa}. Leant and Djex were consulted
for the current source-owned synthesis and provider-selection surface, not to
re-present the verification-boundary review Forge already
contains~\cite{leant,djex,forge-ideas}.

Oink, Spot, and Storm are concrete candidates for scaling search, not dependencies
of the delivered Python code. The mathematics of coupling-based relational
verification and parity solving is established prior work~\cite{oink,coupling}.
All adaptations and proofs below are stated directly so that neither an external
solver's verdict nor a citation is needed to check an individual certificate.
